\documentclass{article}
\usepackage[T1]{fontenc}
\usepackage{lmodern}
\usepackage{graphicx}
\usepackage{amsmath,amssymb}
\usepackage[a4paper,margin=2cm]{geometry}
\usepackage[absolute,overlay]{textpos}
\setlength{\TPHorizModule}{1mm}
\setlength{\TPVertModule}{1mm}
\usepackage[indonesian]{babel}
\usepackage{hyperref}
\setlength{\emergencystretch}{2em}
% unit: Halaman 59

\begin{document}

% === Halaman 59 (llm) ===

\begin{itemize}
    \item Dalam Bahasa Inggris, 10 huruf yang yang paling sering muncul adalah E, T, A, O, I, N, S, H, R, dan D,
    \item Triplet yang paling sering muncul adalah THE. Karena \textbf{LJV} paling sering muncul di dalam cipherteks, maka dari 10 huruf tsb semua kemungkinan kata 3-huruf dibentuk dan kata yang yang cocok untuk \textbf{LJV} adalah \underline{THE}.
    \item Jadi, kita dapat menerka bahwa \textbf{LJV} mungkin adalah \underline{THE}.
    \item Huruf-huruf kunci lainnya dicoba dengan menerka dan menguji coba.
    \item Dari sini kita buat tabel yang memetakan huruf plainteks dengan cipherteks dan huruf-huruf kuncinya (ingatlah bahwa setiap nilai numerik dari huruf kunci menyatakan jumlah pergeseran huruf pada \textit{Caesar cipher}):
\end{itemize}

\end{document}
